\begin{casebox}
\textbf{知识点小案例：日志对象为什么要冷热拆分。}
特例是一条日志记录同时保存 \code{timestamp}、\code{user\_id}、\code{event\_type}、\code{value}、\code{raw\_line} 和 \code{debug\_tags}。聚合热路径只需要 \code{event\_type} 和 \code{value}，但 AoS 布局会把原始文本指针和调试字段一起拖进 cache line。改成热列 \code{event\_type[]}、\code{value[]} 加冷表回查后，热循环读字节减少，错误报告仍能通过 record id 找回原始行。由此推出规律：布局优化不是把结构体换成数组，而是用访问矩阵证明哪些字段同频访问。工程上任何冷热拆分都要同时保留语义、错误身份和回查路径。
\end{casebox}
